\section{\textcolor{Turquoise}{Aufstellen der Hypothesen}} % Statistische Überprüfung der Hypothesen
\label{ch:Ergebnisse:sec:Hypothesen}
%\textcolor{blue}{\textit{einleitende Worte}}\\
Bevor ein Versuch durchgeführt werden kann, werden Hypothesen aufgstellt, um festzulegen welche Ergebnisse von der Untersuchung erwartet werden. Folgende Hypothesen wurden für diesen Versuch aufgestellt:\\

    % -> Usability besser bewerten (Prümper)
    % Aufgaben schneller bearbeitet (Test Zeit)
    %-> Aufgaben verlässlicher beantwortet (Test xo)

% **************************************************************
\noindent
\textbf{1) Wenn die Versuchsperson die Aufgaben mit der neuen Oberfläche bearbeitet, dann kann sie diese schneller beantworten im Vergleich zur alten Oberfläche.} \\

\begin{itemize}
    \item H$_{01}$ = Die Nutzung des neuen Systems ($\mu2$) hat keinen oder keinen positiven Einfluss auf die durchschnittliche Bearbeitungsdauer der gestellten Aufgaben im Vergleich zur Nutzung des alten Systems ($\mu1$); $\mu1 \leq \mu2$.
    \item H$_{11}$ = Die Nutzung des neuen Systems ($\mu2$) beeinflusst die durchschnittliche Bearbeitungsdauer der gestellten Aufgaben positiv im Vergleich zur Nutzung des alten Systems ($\mu1$); $\mu1 > \mu2$.
\end{itemize}

% **************************************************************
\noindent
\textbf{2) Wenn die Versuchsperson die Aufgaben mit der neuen Oberfläche bedient, dann kann sie diese zuverlässiger beantworten als mit der alten Oberfläche} \\

\begin{itemize}
    \item H$_{02}$ = Die Nutzung des neuen Systems ($\mu2$) hat keinen oder keinen positiven Einfluss auf die Fehlerquote der Anworten der gestellten Aufgaben im Vergleich zur Nutzung des alten Systems ($\mu1$); $\mu1 \leq \mu2$.
    \item H$_{12}$ = Die Nutzung des neuen Systems ($\mu2$) beeinflusst die Fehlerquote der Anworten der gestellten Aufgaben positiv im Vergleich zur Nutzung des alten Systems ($\mu1$); $\mu1 > \mu2$.
\end{itemize}

% **************************************************************
\noindent
\textbf{3) Wenn die Versuchsperson die Aufgaben mit der neuen Oberfläche bedient, dann bewertet sie die Usability besser im Vergleich zur alten Oberfläche.}
% -> Usability besser bewerten (Prümper) \\

\begin{itemize}
    \item H$_{03}$ = Die Nutzung des neuen Systems ($\mu2$) hat keinen oder keinen positiven Einfluss auf die Bewertung der Software im Vergleich zur Nutzung des alten Systems ($\mu1$); $\mu1 \leq \mu2$.
    \item H$_{13}$ = Die Nutzung des neuen Systems ($\mu2$) beeinflusst die Bewertung der Software positiv im Vergleich zur Nutzung des alten Systems ($\mu1$); $\mu1 > \mu2$.
\end{itemize}

% **************************************************************
\noindent
Die genannten Hypothesen sind auf gepaarte Daten ausgelegt. Sie sind einseitig, da eine Verbesserung erwartet wird, und die Null-Hypothesen können ab einem Signifikanzniveau von 0,05 wiederlegt werden.
